\documentclass[11pt]{article}
\usepackage[T1]{fontenc}
\usepackage[utf8]{inputenc}
\usepackage[a4paper,margin=25mm]{geometry}
\usepackage{amsmath,amssymb,amsthm,microtype,xcolor}
\usepackage[colorlinks=true,linkcolor=teal!70!black,urlcolor=teal!70!black,citecolor=teal!70!black]{hyperref}
\newtheorem{proposition}{Proposition}
\setlength{\parindent}{0pt}
\setlength{\parskip}{5pt}
\title{Ordered Virtual Scales and the Leading Soft-Photon Logarithm}
\author{\normalsize Deep Bhattacharjee, James Montgomery Lockwood, Tristan Romanov,\\
\normalsize Ikailo John Sessford, Ushashi Bhattacharya, Sanjeevan Singha Roy}
\date{Working research draft, 15 November 2026}

\begin{document}
\maketitle
\begin{center}\small\color{teal!70!black}
One unranked prospective author list supplied by the project lead.\\
Coauthor review and approval remain pending.
\end{center}
\begin{abstract}
Ordered virtual scales give a simplex of logarithmic integration variables.
We prove that an integrable bound in the adjacent or endpoint gaps
places the integrated discrepancy at one lower logarithmic power.
The estimate is sharp for a nonnegative single-gap correction.
Bounded primitives give another sufficient condition for signed
single-gap corrections. For the celestial single-soft-photon contact
term in massless scalar QED, this identifies an estimate that would
suffice to recover the two-loop leading coefficient from a factorized
separated-scale kernel. The interacting ordered kernel remains to be
derived.
\end{abstract}

\section{A sufficient estimate for ordered logarithmic scales}
\label{sec:orderedscales}

Choi, Kadhe, and Puhm identify the persistence of the logarithmic
soft-factor tower beyond one loop as an open quantum test
\cite[Sec.~6]{longrange}. For massless scalar QED, the first test is
the coefficient of $\omega\log^2(\omega/\Lambda)$ in the
symmetry-governed single-soft-photon insertion, with hard scale
$\Lambda$. In their positive-helicity convention, define the
tree prefactor $S_{-1,i}=\sqrt{2}eQ_i/(z-z_i)$ and write
\[
 B_i=-\gamma_{\mathrm{em}}\eta_i
       \sum_{j\ne i}\frac{(\bar z-\bar z_j)Q_iQ_j}
                          {\bar z_i-\bar z_j}\,
         e^{-\partial_{\Delta_i}} .
\]
Here $z_i$ are distinct hard punctures, $Q_i$ and $\eta_i$ are
the charges and in/out signs in their convention,
$\gamma_{\mathrm{em}}=e^2/(4\pi^2)$, and the exponential
shifts the hard weight $\Delta_i$. Equation~(4.24) of
\cite{longrange} uses the normalized Mellin coefficient
in their Eq.~(4.11). In this normalization the energy
coefficient is preserved, since
\[
 \int_0^\Lambda \omega^{\Delta-1}\omega\log^2(\omega/\Lambda)\,d\omega
 =\frac{2\Lambda^{\Delta+1}}{(\Delta+1)^3}
 \qquad(\Re\Delta>-1)
\]
and multiplication by $(\Delta+1)^3/2$ extracts its leading
coefficient at $\Delta=-1$. The predicted contribution is
$S_{-1,i}B_i^2H/2$ on a fixed hard correlator $H$. To test
the contact coefficient at
puncture $i$, pair $\bar\partial^3$ of this term with a smooth
test supported near $z_i$ and away from the other hard punctures,
as motivated by the descendant in Eq.~(4.25). This projects the
product after composition. The following estimate gives
conditions on an independently calculated kernel sufficient
to obtain the $1/2$ in the projected leading coefficient.

\begin{proposition}[Ordered-scale leading coefficient]
\label{prop:orderedscales}
Let $X$ be a Banach space and $m\ge2$. For $T>0$ use the ordered domain
\[
 S_m(T)=\{(t_1,\ldots,t_m):0<t_1<\cdots<t_m<T\}.
\]
Suppose
\[
 F_m(T)=\int_{S_m(T)}K_T(t)\,dt+E_T,\qquad
 \|E_T\|\le C T^{m-1},
\]
where $K_T$ is Bochner integrable. Define $d_0=t_1$,
$d_j=t_{j+1}-t_j$ for $1\le j<m$, and $d_m=T-t_m$.
Suppose there are $A_m\in X$ and nonnegative
$w_0,\ldots,w_m\in L^1([0,\infty))$, independent of $T$, such that
\[
 \|K_T(t)-A_m\|\le\sum_{j=0}^m w_j(d_j)
 \quad\text{for almost every }t\in S_m(T).
\]
Then
\begin{equation}
 \left\|F_m(T)-\frac{T^m}{m!}A_m\right\|
 \le T^{m-1}\left(C+\frac1{(m-1)!}\sum_{j=0}^m\|w_j\|_1\right).
 \label{eq:orderedscale-error}
\end{equation}
\end{proposition}
\begin{proof}
The change from the ordered $t_j$ to the independent gaps
$d_0,\ldots,d_{m-1}$ has unit Jacobian and gives the simplex
of nonnegative gaps with sum at most $T$. The remaining gap
is $d_m=T-\sum_{j=0}^{m-1}d_j$. Any gap can be chosen as the
dependent coordinate, since the corresponding affine
substitution has Jacobian of absolute value one. At fixed
$d_j=u$, the other independent coordinates form a simplex
of dimension $m-1$ with volume $(T-u)^{m-1}/(m-1)!$.
Consequently
\begin{equation}
 \int_{S_m(T)}w_j(d_j)\,dt=
 \frac1{(m-1)!}\int_0^T (T-u)^{m-1}w_j(u)\,du
 \le\frac{T^{m-1}}{(m-1)!}\|w_j\|_1 .
 \label{eq:orderedscale-gap}
\end{equation}
Since $S_m(T)$ has volume $T^m/m!$, the constant part
integrates to $T^m A_m/m!$. Apply the integral norm
inequality to the difference and include the bound on $E_T$.
\end{proof}

The error power in \eqref{eq:orderedscale-error} is sharp.
With $E_T=0$ and $K_T-A_m=e^{-d_j}v$ for a unit vector $v$,
the integrated norm is asymptotic to $T^{m-1}/(m-1)!$ by
\eqref{eq:orderedscale-gap} and dominated convergence.
A slowly decaying choice $(1+d_j)^{-1}v$ instead produces
\[
 \frac{T^{m-1}}{(m-1)!}\big(\log(1+T)+O(1)\big).
\]
Indeed, the difference between the scaled integral in
\eqref{eq:orderedscale-gap} and $\int_0^T (1+u)^{-1}\,du$
is bounded using $0\le1-(1-u/T)^{m-1}\le (m-1)u/T$.
Thus pointwise decay at large gap does not guarantee the
error rate in \eqref{eq:orderedscale-error}. For $m=2$, a
constant additional kernel $v$ would give a contribution
$T^2v/2$, even with all one-loop data held fixed.

For the QED test, take $T=\log(\Lambda/\omega)$ and write
virtual scales as $\kappa_j=\omega e^{t_j}$, so
$d\kappa_j/\kappa_j=dt_j$. To apply
Proposition~\ref{prop:orderedscales}, one must extract
$F_2$ after dividing by the single soft-energy power
$\omega$ and projecting in a fixed regulator and
soft-limit prescription. Combining all contributing
orderings and normalizations, controlling collinear and
overlap subtractions, and bounding omitted regions by
$O(T)$ are part of deriving the ordered representation.
The well-separated kernel must then match the projected
$S_{-1,i}B_i^2H$, with its deviation controlled also at
comparable scales and the two endpoints. A derivation of
ordinary virtual infrared exponentiation, such as the
gravity calculation of Naculich and Schnitzer
\cite{naculichschnitzer}, concerns a virtual divergence
with different kinematics. The ordered representation
and kernel estimate for the radiative QED coefficient
remain to be established.


\section{Absolute gap bounds and cancellation}
For a nonnegative scalar function $w$ integrable on finite intervals,
the contribution from a single gap in \eqref{eq:orderedscale-gap}
satisfies
\[
 \frac{(T/2)^{m-1}}{(m-1)!}\int_0^{T/2} w(u)\,du
 \le \int_{S_m(T)} w(d_j)\,dt
 \le \frac{T^{m-1}}{(m-1)!}\int_0^T w(u)\,du .
\]
The lower bound follows by restricting the one-gap integral to
$0\le u\le T/2$. Hence this contribution is $O(T^{m-1})$ as
$T\to\infty$ exactly when $w\in L^1([0,\infty))$.

For a signed or complex scalar profile $a$ that is integrable on finite
intervals, let $A(u)=\int_0^u a(v)\,dv$. If $A$ is bounded, then for each gap
\[
 \int_{S_m(T)}a(d_j)\,dt
 =\frac1{(m-2)!}\int_0^T(T-u)^{m-2}A(u)\,du .
\]
The same simplex formula applies by linearity on finite intervals.
Integration by parts has zero boundary terms because $A(0)=0$ and
$m\ge2$. It gives an absolute bound of
$\|A\|_\infty T^{m-1}/(m-1)!$. For instance,
$a(u)=\sin(u)$ has bounded primitive $1-\cos(u)$ and nonintegrable
absolute value. The example shows that a signed remainder can satisfy
the desired integrated bound through cancellation. This is relevant
after summing diagrams and subtraction terms, where the remainder can
be signed or complex.

\paragraph{Authorship and preparation.}
The unranked prospective byline matches the companion manuscript.
Coauthor review, contribution statements, and approval remain pending.
An AI coding assistant assisted the derivation, literature search,
and drafting.

\begin{thebibliography}{9}
\bibitem{longrange}
S. Choi, A. Kadhe, and A. Puhm, \emph{Long-Range Interactions in Celestial CFT},
arXiv:2601.05951v1 (2026), \url{https://arxiv.org/abs/2601.05951v1}.
\bibitem{naculichschnitzer}
S. G. Naculich and H. J. Schnitzer,
\emph{Eikonal Methods Applied to Gravitational Scattering Amplitudes},
JHEP 05 (2011) 087, \url{https://arxiv.org/abs/1101.1524}.
\end{thebibliography}
\end{document}
